\documentclass[parskip=full,11pt,twoside]{scrartcl}
\usepackage[utf8]{inputenc}

\begin{titlepage}

\subject{Pflichtenheft}
\title{\textbf{$\lambda$}urora}
\subtitle{The Lambda Calculus IDE}


\author{Alexander von Heyden, Julia Patrusheva\\
 Max Nowak, Nikolai Polley\\
 Randy Seng, Younis Bensalah}

\newpage
\end{titlepage}

% section numbers in margins:
\renewcommand\sectionlinesformat[4]{\makebox[0pt][r]{#3}#4}

% header & footer
\usepackage{scrlayer-scrpage}
\lofoot{\today}
\refoot{\today}
\pagestyle{scrheadings}

\usepackage[sfdefault,light]{roboto}
\usepackage[T1]{fontenc}
\usepackage[german]{babel}
\usepackage[yyyymmdd]{datetime} % must be after babel
\renewcommand{\dateseparator}{-} % ISO8601 date format
\usepackage{hyperref}
\usepackage[nameinlink]{cleveref}
\crefname{figure}{Abb}{Abb}
\usepackage[section]{placeins}
\usepackage{xcolor}
\usepackage{graphicx}
\hypersetup{
	pdftitle={Pflichtenheft},
	bookmarks=true,
}
\usepackage{csquotes}

\usepackage{amsmath} % for $\text{}$
\newcommand\urlpart[2]{$\underbrace{\text{\texttt{#1}}}_{\text{#2}}$}
\usepackage{color}

\usepackage{pflichtenheft}




\begin{document}

\maketitle
\newpage
\tableofcontents

\newpage
\section{Einleitung}
Der Lehrstuhl für Programmierparadigmen am IPD Snelting ist seit Jahren dafür zuständig, Besuchern der gleichnamigen Vorlesung des $\lambda$-Kalkül zu vermitteln. Leider hat es sich in den vergangenen Jahren immer wieder gezeigt, dass eine erhebliche Anzahl von Studenten mit den grundlegenden Eigenschaften des $\lambda$-Kalküls Probleme haben. Hier versucht Aurora anzusetzen. Mit Hilfe einer einfach strukturierten Web App, die vom Design her an gängige IDEs angelehnt ist, soll dem Benutzer die Einführung in das $\lambda$-Kalkül erleichert werden. Hierfür wird dem Besucher eine Eingabemaske zur Verfügung gestellt, in die er eigene Terme eingeben kann. Schritt für Schritt kann er nun die Auswertung bis zum finalen Ergebnis verfolgen.


\newpage
\section{Zielsetzung}
% Diese Section sollte kurz und knapp "für Manager" sein
% und auf eine Seite passen.


\subsection{Musskriterien}

\criterium{Eingabe}{crt:text-editor}
Es gibt ein Eingabefeld um $\lambda$-Terme einzugeben.

\criterium{Input editieren}{crt:not-editable}
Der Input kann nur bearbeitet werden, solange das Programm nicht rechnet.

\criterium{Auswertung starten}{crt:evaluate-term}
Der Benutzer kann die Auswertung des eingegebenen $\lambda$-Terms starten. Im Ausgabefenster erscheint das Ergebnis in Normalform. 

\criterium{Auswertung stoppen}{crt:stop-feature}
Der Benutzer kann die Auswertung des $\lambda$-Terms jederzeit stoppen.

\criterium{Zwischenschritte}{crt:intermediate-step-output}
Nach der Auswertung des $\lambda$-Terms können Zwischenergebnisse ausgegeben werden.

\criterium{Auswertungsstrategie}{crt:default-eval-strategy}
Die Auswertungsstrategie bei $\beta$-Reduktion von $\lambda$-Termen ist die Normalreihenfolge. Weitere Auswertungsstrategien sind in den Wunschkriterien.


\newpage
\subsection{Wunschkriterien}

\criteriumOptional{Zeilennummer}{crt:linenumbers}
Neben dem Eingabefeld soll ausgeblasst eine Zeilennumerierung dargestellt werden um eine einfache Zeilenzuteilung zu ermöglichen.

\criteriumOptional{Kommentare einfügen}{crt:comments}
Man kann eine Zeile in der Eingabe auskommentieren. Diese Zeile wird dann vom Parser übersprungen

\criteriumOptional{Autovervollständigung von Klammern}{crt:autobracket}
Wenn der Benutzer eine Klammer öffnet \enquote{(} dann schließt das Programm automatisch diese Klammer mit \enquote{)}

\criteriumOptional{Syntax-Highlighting}{crt:syntax-highlighting}
Schlüsselwörter aus den Bibliotheken und $\lambda$s werden farblich hervorgehoben. Beim hovern über einer Klammer wird das entsprechende Klammerpaar farblich hervorgehoben.

\criteriumOptional{Erweitertes Syntax-Highlighting}{crt:syntax-highlighting-premium}
Der nächste auszuwertende Redex wird farblich hinterlegt, im nächsten Schritt in einem helleren Ton derselben Farbe.

\criteriumOptional{Standardbibliothek}{crt:std-library}
Bibliothek, in der vordefinierte Funktionen und Church-Zahlen verfügbar sind. Durch Schlüsselwörter kann der Benutzer diese direkt in seinen Code einbinden.

\criteriumOptional{Benutzerbibliothek}{crt:user-library}
Bibliothek, in der der Benutzer eigene Funktionen definieren und speichern kann. Durch Schlüsselwörter kann der Benutzer diese direkt in seinen Code einbinden.

\criteriumOptional{Weitere Auswertungsstrategien}{crt:opt-eval-strategy}
Zusätzliche Auswertungsstrategien sind \enquote{Call by Name} und \enquote{Call by Value}.

\criteriumOptional{Benutzerdefinierte Auswahl}{crt:own-eval-strategy}
Eigene Auswahl bei der $\beta$-Reduzierung eines Redex.

\criteriumOptional{Pause-Button}{crt:pausebutton}
Der Pause-Button unterbricht die laufende Berechnung
 
\criteriumOptional{Share-Button}{crt:share-feature}
Ein Share-Button ermöglicht das Erstellen von short-URLs sowie das Kopieren dieser in den Zwischenspeicher.

\criteriumOptional{LaTeX export}{crt:latex-export}
$\lambda$-Terme aus dem Text-Editor oder der Schrittliste können in LaTeX-Snippet exportieren werden. Bis zu einer gewissen größe ist es auch möglich, den kompletten Rechenweg mit zu exportieren.

\criteriumOptional{$\lambda$-Term vom Text-Editor oder Schrittliste in LaTeX-Snippet exportieren}{crt:latex-export}

\criteriumOptional{LaTeX export}{crt:latex-export}
$\lambda$-Terme aus dem Text-Editor oder der Schrittliste können in LaTeX-Snippet exportieren werden. Bis zu einer gewissen Größe ist es auch möglich, den kompletten Rechenweg mit zu exportieren.


\criteriumOptional{\enquote{Night Mode}}{crt:night-mode}
Die Funktion \enquote{Night Mode} ändert das äußere Erscheinungsbild der Applikation zu einem dunkleren Farbschema.

\criteriumOptional{Schleifenerkennung}{crt:loop-detection}
Das Programm soll einfache Schleifen wie z.B. ($\lambda$x. x x) ($\lambda$x. x x) erkennen. Insbesondere kann es aber weiterhin vorkommen, dass komplexere Schleifen nicht erkannt werden.

\criteriumOptional{$\lambda$urora Tutorial}{crt:tutorial}
Das Tutorial erklärt und zeigt dem Benutzer, wie man die Applikation $\lambda$urora nutzt.

\criteriumOptional{De Bruijn Indizes}{crt:debruijn-output}
Ausgabe des Programmes kann wahlweise mit De Bruijn Indizes erfolgen.

\criteriumOptional{Lokale Sessions}{crt:local-sessions}
Die Applikation merkt sich im localStorage des Webbrowsers die Historie von vorherigen Eingaben. Der Benutzer kann außerdem diese Historie verändern.

\criteriumOptional{Tastatur-Shortcuts}{crt:shortcuts}
$\lambda$urora stellt Tastatur-Shortcuts zum vereinfachten Ausführen von Features bereit.

\textbf{Verfügbare Tastatur-Shortcuts}

\begin{tabular}{|r|l|}
    
    \hline 
    shortcut & Beschreibung \\ \hline
    \textbackslash & $\lambda$ \\ \hline
    Ctrl + $\hookleftarrow$ & Run \\ \hline
    
\end{tabular}


% Hilft den Entwicklern beim Debuggen.













\subsection{Abgrenzungskriterien}

\criteriumNot{Berechnungen laufen Client-seitig.}{crt:client-side}
Der Client führt das Programm aus. Alle Funktionen die etwas berechnen oder speichern müssen werden auf der Hardware des Clients ausgeführt.

\criteriumNot{Keine Server-Komponente.}{crt:no-server-side}
Der Server ist nur für die Verbreitung des Programmcodes zuständig. Er selber speichert keinerlei Daten oder führt Aktion für das Programm aus.

\criteriumNot{Kein typisiertes $\lambda$-Kalkül.}{crt:no-typed-calculus}
Das Programm kann nur $\beta$-Reduktion des untypisierten $\lambda$-Kalküls. Auch eine einfache Erweiterungsschnittstelle, damit man z.B. das typisierte $\lambda$-Kalkül einfügen könnte ist nicht geplant und wird nicht bei Entwicklung des Programms in Betracht gezogen

\criteriumNot{Applikation bietet keine Projekt-Unterstützung.}{crt:no-projects}
Das Programm unterstützt keine Klassen oder anderweitig mehr als eine Eingabe oder mehrere Ausgabe.
Das Programm ist so konzipiert das ein Benutzer es nur alleine bedienen kann und gegebenenfalls mit der Share-Funktion einnem anderem die Ergebnisse des Programmes mitteilen kann. Ein Versionskontrolle oder generelles Zusammenarbeiten über zwei Clients mit dem gleichen Programm ist nicht möglich.



\pagebreak
%%%%%%%%%%%%%%
\section{Produkteinsatz}

Wenn ein Interessierter das Programm benutzen will, dann benötigt er eine Internetanbindung. Er benötigt  außerdem einen PC mit folgenden Spezifikationen: Wenn er Windows benutzt, dann Windows 7 oder höher, wenn er Mac benutzt dann OS X Mavericks 10.9 oder höher oder wenn er Linux benutzt Ubuntu 14.04 oder höher, Debian 8 oder höher, oder Fedora24 oder höher.
Er braucht zusätzlich einen SSE2-fähigen Intel Pentium 4-Prozessor oder besser.

\subsection{Anwendungsbereiche}

Der Anwendungsbereich umfasst die Studenten und wissenschaftliche Mitarbeiter, die sich mit dem $\lambda$-Kalkül beschäftigen. Aurora erlaubt dem Benutzer $\lambda$-Terme zu berechnen und sich dabei einzelne Schritte der $\beta$-Reduktion anzeigen zu lassen. Durch vorhandene Shortcuts und eine mit grundlegenden Funktionen ausgestattete Standardbibliothek soll die Benutzung der Aurora Web App den Benutzer an eine normale IDE erinnern, sodass schneller das gewünschte Ergebnis erreicht werden kann.

\subsection{Zielgruppen}

Zu der Zielgruppe gehören vor allem Informatik-Studenten und wissenschaftliche Mitarbeiter, die sich mit dem $\lambda$-Kalkül beschäftigen. Hierzu gehört z.B. das schnelle Auswerten eines vorhandenen $\lambda$ Terms, oder auch das schrittweise durchgehen einer Auswertung zum selbstständigen Üben. Auch anderweitig Interessierte sollen das Programm benutzen können.
Die Plattform ist primär für Einzelpersonen gedacht, eine Gruppenarbeit mit dem Programm ist nur eingeschränkt möglich. Die Aurora Web App soll als Werkzeug genutzt werden, und ihre Zielgruppen somit während des Studiums, wissenschaftlichen Arbeitens oder im privaten Umfeld unterstützen.


Für die Benutzung der Aurora Web App werden Basiskenntnisse der Internetnutzung vorausgesetzt. Auch vorausgesetzt werden außerdem Kenntnisse über das untypisierten $\lambda$-Kalkül.

Die Basissprache der Web App ist Englisch, daher sind Englisch Grundkenntnisse von Vorteil.

\subsection{Betriebsbedingungen}
Bezüglich den Betriebsbedingungen, soll sich Aurora nicht wesentlich von anderen Web Apps unterscheiden. Die maximale Anzahl der Benutzer von Aurora ist unbegrenzt bzw. vom Webserver abhängig. (Die Betriebsdauer soll 24 Stunden betragen.) 


\section{Produktumgebung}



%%%%%%%%%%%
\section{Funktionale Anforderungen}

\subsection{Eingabe}
\functionality{Eingabefeld}{fnc:input}\fulfills{crt:text-editor}\fulfills{crt:not-editable}
Der Benutzer kann $\lambda$-Terme über das Eingabefeld mit einer Tastatur eingeben. Die Zeichen erscheinen in Echtzeit im Eingabefeld.

\functionality{Eingabefeld bearbeiten}{fnc:editor}\fulfills{crt:text-editor}\fulfills{crt:not-editable}
Der Benutzer kann den Inhalt das Eingabefelds durch gängige Methoden bearbeiten.
Er kann mit der Maus Text auswählen, mit \enquote{Delete} Zeichen löschen usw.

\functionality{Run Eingabefeld}{fnc:run-input}\fulfills{crt:evaluate-term}
Der \enquote{Run-Button} startet die  \enquote{Run}-Funktion und diese veranlasst den Parser dazu, dass er die $\lambda$-Terme über das Eingabefeld einliest.

\functionality{Shortcuts}{fnc:output}
\functionality{Auto-complete}{fnc:output}
\functionality{Highlightning Lambda}{fnc:output}
\functionality{Highlightning Text}{fnc:output}

\functionality{Brackets Auto-complete}{fnc:bracket}
\fulfills{crt:autobracket}
Wenn der Benutzer eine Klammer öffnet \enquote{(} dann erkennt das Programm in Echtzeit, dass eine Klammer geöffnet wurde und fügt automatisch eine schließende Klammer \enquote{)} hinzu

\functionality{Standardbibliothek}{fnc:stdlib}\fulfills{crt:std-library}
Der Benutzer kann Funktionen aus der Standardbibliothek im Eingabefeld verwenden.

\functionality{Funktionen der Benutzerbibliothek}{fnc:user-vs-std}\fulfills{crt:user-library}
Funktionen aus der Benutzerbibliothek können analog zu Funktionen der Standardbibliothek verwendet werden. Es unterscheidet sich für die Eingabe nicht, ob eine Funktion in der Standard oder in der Benutzerbibliothek ist.


\functionality{Äquivalenz von Funktionen und deren Definition}{fnc:func-lambda-equivalent}\fulfills{crt:user-library}\fulfills{crt:std-library}
Die Eingabe von definierten Funktionen aus der Standard- und Benutzerbibliothek ist semantisch äquivalent zur Eingabe deren entsprechenden Definition als $\lambda$-Term.

\functionality{Auswahl der Auswertungsstrategie}{fnc:eval-order}\fulfills{crt:opt-eval-strategy}
Der Benutzer kann die Auswertungsstrategie der $\lambda$-Terme auswählen.

\functionality{Auswertungsstrategien}{fnc:eval-order-options}\fulfills{crt:opt-eval-strategy}\fulfills{crt:own-eval-strategy}
Der Benutzer kann zwischen 3 Standard Auswertungsstrategien auswählen. Sie sind \enquote{Normalreihenfolge}, \enquote{Call-by-Name} und \enquote{Call-by-Value}.
Er kann außerdem auf einen Redex drücken der dann ausgewertet wird.

\functionality{Der Benutzer kann selber Redex auswählen, der als nächstes reduziert wird}{fnc:choosingredex}\fulfills{crt:own-eval-strategy}
Wenn der Benutzer sich die Steps anzeigen lässt kann er auf  einen Redex drücken und der wird als nächstes ausgewertet. 
\functionality{Die Eingabe kann nur geändert werden, falls keine Berechnung läuft.}{fnc:no-edit-while-running}\fulfills{crt:not-editable}
Der Benutzer kann das Eingabefeld nur genau dann editieren, wenn keine Berechnung am laufen ist.

\functionality{Ausgabefeld nicht editierbar}{fnc:output-readonly}\fulfills{crt:not-editable}
Der Inhalt des Ausgabefelds kann nicht bearbeitet werden.

\functionality{Kommentare hinzufügen}{fnc:commentyourcode}\fulfills{crt:comments}
Der Benutzer kann mit \enquote{\#} die ganze Zeile nach dem Zeichen auskommentieren
Der Parser erkennt das Zeichen und springt zur nächsten Zeile.

\subsection{Darstellungen}
\functionality{Night Mode}{fnc:night-mode}\fulfills{crt:night-mode}
Das äußere Erscheinungsbild der Applikation kann über die \enquote{Night mode}-Funktion zu einem dunklen Farbschema geändert werden. Dies kann über das Hamburger-Menü ein- und ausgestellt werden.

\functionality{Line Numbers}{fnc:ln}\fulfills{crt:linenumbers}
Links neben der Eingabe des Benutzers sind Zahlen, die jeweils die jetzige Zeilennummer angeben. Sie sind nicht editierbar oder anderweitig änderbar. Dies soll mit der Orientierung der Eingabe helfen.


\functionality{Tutorial zur Web App}{fnc:tutorial}
Wenn der Benutzer auf einen Hilf-Button drückt dann wird ein kleines Tutorial erscheinen, dass ihm die Benutzung des Programmes erklärt und auf Features hinweist.

\functionality{Ausgabefeld}{fnc:output}\fulfills{crt:evaluate-term}
Falls die Berechnung terminiert, steht das Ergebnis im Ausgabefeld.

\functionality{Nach Wunsch der User werden die $\beta$-Reduktionsschritte angezeigt}{fnc:stepbystep}\fulfills{crt:intermediate-step-output}
Dem Benutzer wird in der GUI nach einer erfolgreichen Berechnung 3 Punkte angezeigt. Wenn er auf diese klickt werden ihm alle $\beta$-Reduktions Schritte angezeigt die durchgeführt wurden um auf das Ergebnis zu kommen


\functionality{Highlighting des nächsten auszuwertenden Redex}{fnc:highlight-next}\fulfills{crt:syntax-highlighting-premium}
In jedem Schritt wird im Schrittfeld nächste auszuwertende Redex farblich hinterlegt.

\functionality{Highlighting des zuletzt ausgewerteten Redex}{fnc:highlight-prev}\fulfills{crt:syntax-highlighting-premium}
In jedem Schritt wird der zuletzt eingesetzte Term farblich hinterlegt.

\functionality{Finales Ergebnis wird angezeigt}{fnc:output}
\textcolor{green}{HIER FEHLT BESCHREIBUNG}


\functionality{Sprache ist auswählbar}{fnc:language-selection}
Der Benutzer kann aus einem Satz von menschlichen Sprachen, die Sprache der Applikation einstellen.

\fulfills{crt:default-language}
\fulfills{crt:opt-language}

\functionality{GUI wird korrekt auf Google Chrome Web Browser angezeigt}{fnc:output}
\textcolor{green}{HIER FEHLT BESCHREIBUNG}

\functionality{Die Web App wird korrekt in Google Chrome Web Browser funktionieren}{fnc:output}
\textcolor{green}{HIER FEHLT BESCHREIBUNG}

\functionality{Stop-Button}{fnc:stop-exec} \fulfills{crt:stop-feature}
Der Benutzer kann eine laufende Berechnung über einen \enquote{Stop}-Button vorzeitig beenden. Alle Berechnungsschritte werden verworfen und nichts aus der bisherigen Berechnung wird gespeichert.

\functionality{De Bruijin Indizes}{fnc:debruijn-output}
Ausgabe der $\lambda$-Ausdrücke mit De Bruijn Indizes.


\functionality{Mobile Darstellung ist im begrenzten Modus möglich}{fnc:output}
\textcolor{green}{HIER FEHLT BESCHREIBUNG}

\subsection{Hintergrundablauf}
\functionality{Run}{fnc:run}\fulfills{crt:evaluate-term}
Der Benutzer kann die Berechnung des $\lambda$-Terms über den \enquote{Run}-Button starten.

\functionality{Eingabe wird geparsed}{fnc:output}
\textcolor{green}{HIER FEHLT BESCHREIBUNG}

\functionality{$\alpha$-Konversion mit De Bruijn Indices}{fnc:output}
\textcolor{green}{HIER FEHLT BESCHREIBUNG}

\functionality{Normalreihenfolge}{fnc:eval-order-normal}\fulfills{crt:default-eval-strategy}
$\lambda$-Terme werden standardmäßig in Normalreihenfolge ausgewertet.
Das heißt, dass immer der linkeste Redex ausgewertet wird.

\functionality{Ermittlung der Eingestellten Auswertungsstrategie}{fnc:output}
\textcolor{green}{HIER FEHLT BESCHREIBUNG}

\functionality{Offensichtliche Schleifen erkennung}{fnc:loop-test}
Eine laufende Berechnung wird automatisch auf Schleifen überprüft. Wenn eine Schleife gefunden wurde warnt das Programm davor.

\functionality{$\beta$-Reduktion}{fnc:output}
\functionality{Das Program erkennt wenn keine weitere $\beta$-Reduktion möglich ist}{fnc:output}

\functionality{Pause-Button} {fnc:pausebut} \fulfills{crt:pausebutton}
Der Benutzer kann auf den Pause-Button drücken um die laufende Berechnung zu unterbrechen und sich z.B. die bisherigen Schritte anzeigen lassen. Er kann allerdings im Gegensatz zum Stop-Button die Berechnung weiterlaufen lassen von dem Punkt aus, an dem sie gestoppt hat. 
\functionality{No Timeout}{fnc:run-no-timeout}\fulfills{crt:infiniteloop}
Ein vom Benutzer eingegebenes Programm kann potentiell unendlich laufen.

\subsection{Speichern}

\functionality{Funktionen der Standardbibliothek}{fnc:func-name-def}\fulfills{crt:std-library}
Funktionen aus der Standardbibliothek haben einen Namen und eine Definition ($\lambda$-Term).

\functionality{Definition einer Funktion}{fnc:func-def}\fulfills{crt:std-library}\fulfills{crt:user-library}
Eine definierte Funktion der Standardbibliothek oder Benutzerbibliothek entspricht genau einem $\lambda$-Term.

\functionality{Benutzerbibliothek}{fnc:userlib}\fulfills{crt:user-library}
Der Benutzer kann Funktionen zur Benutzerbibliothek hinzufügen.


\functionality{Share-Button}{fnc:share-button}
Der Benutzer kann mithilfe des Share-Buttons seinen Input in einer short-URL zum Teilen speichern. Wenn ein anderer Benutzer diesen Link in seinen Browser einfügt, dann öffnet sich das Programm mit dem gleichen Input, wie derjenige der den Link geschickt hat.

\functionality{Session und Benutzer Bibliothek Speichern}{fnc:output}
\textcolor{green}{HIER FEHLT BESCHREIBUNG}

\functionality{Berechnung in LaTeX schreiweise umwandeln und zum Copy-to-Clipboard hinzufügen}{fnc:output}
\textcolor{green}{HIER FEHLT BESCHREIBUNG}

\functionality{Stop-Button}{fnc:but}
Wenn der Stop-button betätigt wird muss er ohne große Latenz und und ohne einfrieren des Programmes das Berechnen beenden können.


\functionality{Anschaufunktion auf mobilen Geräten}{fnc:gui-mobile-devices}
Die Funktionalität der Aurora Web App ist auf mobilen Geräten gegenüber der Desktop Version eingeschränkt.

\subsubsection{History}
\functionality{Historie}{fnc:localstorage-autosave}\fulfills{crt:local-sessions}
Eine ausgeführte Benutzereingabe, wird automatisch in die Historie eingefügt.

\functionality{Permanenz der Historie}{fnc:localstorage-permanent}\fulfills{crt:local-sessions}
Einträge in der Historie werden nach schließen des Browsers nicht gelöscht.

\functionality{Wiederherstellen der Historie}{fnc:localstorage-restore}\fulfills{crt:local-sessions}
Der Benutzer kann Eingaben aus der Historie wiederherstellen.

\functionality{Löschen der Historie}{fnc:localstorage-flush}\fulfills{crt:local-sessions}
Der Benutzer kann die Einträge in der Historie löschen.

\functionality{Benutzerbibliothek in Historie}{fnc:localstorage-userlib}\fulfills{crt:local-sessions}
Eine vom Benutzer definierte Funktion wird automatisch in die Historie eingefügt.
%\functionality{Login-Möglichkeit auf Homepage}{fnc:login}
%\fulfills{crt:login}
%\fulfills{crt:github}

%Auf der Homepage \texttt{http://atu.rl/} sieht ein Besucher
%einen \enquote{Login via Facebook} Knopf.
%Weitere Knöpfe wie \enquote{Login via Github} sind möglich.
%Siehe \cref{fig:homepage}.



%%%%%%%%%%%
\section{Nicht-Funktionale Anforderungen}

\nonFunctionality{Benutzbarkeit}{nfc:usability}
Das Programm soll auch ohne dem Tutorial verständlich bedienbar sein.

\nonFunctionality{Aussehen}{nfc:appearence}
Die Grafische Oberfläche soll schlicht aber modern wirken. 

\nonFunctionality{Geschwindigkeit}{nfc:latency}
Das Programm soll sich schnell und responsiv anfühlen. Die Aktion die ein Button ausführt soll schnell erfolgen. 

\nonFunctionality{Programmsprache}{crt:default-language}
Die Standard Sprache des Programmes ist Englisch.

\nonFunctionality{Erweiterte Sprachenauswahl}{nfc:opt-language}
Das Programm ist durch außerdem in den Sprachen Deutsch, Russisch, Französisch verfügbar.

\nonFunctionality{Browserkompatibilität}{nfc:default-browser}
Das Programm funktioniert auf dem Webbrowser \enquote{Google Chrome} 62.

\nonFunctionality{Erweiterte Browserkompatibilität}{nfc:advances-default-browser}
Das Programm funktioniert auf aktuellen Webbrowsern, die Javascript unterstützen.

\nonFunctionality{Unterstützde Sprachen}{nfc:supported-languages}
Die Programmarchitektur soll einfaches Hinzufügen neuer, menschlicher Sprachen unterstützen.

\nonFunctionality{Hash-Aufwand}{nfc:hash-time}
Die Berechnung der Hashwerte ist nicht mehr als 10\% der Gesamtrechenzeit in Anspruch.

\nonFunctionality{Standardbibliothek}{nfc:standard-library-content}
Die Standardbibliothek kann bis zu 20 vordefinierte Funktionen beinhalten.

\nonFunctionality{Benutzerbibliothek}{nfc:user-library-content}
Die Benutzerbibliothek kann bis zu 10 selbstdefinierte Funktionen beinhalten.

\nonFunctionality{Begrenzte Zwischenergebnisse}{nfc:limit-intermediate-results}
Die Anzahl der sichtbaren Zwischenergebnisse ist auf XXX beschränkt.

\nonFunctionality{Begrenzter LaTeX export}{nfc:limited-latex-export}
Der LaTeX-Export für die komplette Rechnung kann nur ausgeführt werden, falls weniger als 20 Zwischenergebnisse vorliegen.



%%%%%%%%%%%
\section{Tests}


%\tests{fnc:impressum-link}
%\tests{fnc:datenschutz-link}
\teststep {Das Programm liegt auf einem Server}
{Der Benutzer tippt den link des Programms \texttt{aurora.kit.edu} in den Browser \enquote {Google Chrome} }
{ Die Webseite wird geöffnet und wird korrekt dargestellt je nach Browserauflösung}

\teststep {Das Programm ist geöffnet und es ist im Default-State}
{ Der Benutzer gibt in den Editor \enquote {($\lambda$ x.x) z } ein. Dannach drückt er auf den Button \enquote {Run}}
{ Das Programm führt eine $\beta$-Reduktion aus und gibt als Ausgabe \enquote {z} aus. Es wird auch angezeigt, dass es nur eine $\beta$-Reduktion durchgeführt hat.}

\teststep {Das Programm ist geöffnet und es ist im Default-State}
{Der Benutzer gibt den Term \enquote {$\lambda$ x.x} in das Eingabefeld ein und drückt auf Run}
{Das Programm kann keine $\beta$-Reduktion durchführen und gibt als Ausgabe $\lambda$ x.x}

\teststep {Das Programm ist geöffnet und findet sich im Default-State. An der rechten Seite ist die Standardfunktionsbibliothek implementiert }
{Der Benutzer gibt in den Editor \enquote {plus 2 2} ein und drückt auf den Buttorn \enquote {Run}}
{ Das Programm erkennt das Wort \enquote {plus} aus der Standardbibliothek und ersetzt es durch \enquote {$\lambda$n.$\lambda$m.$\lambda$s.$\lambda$z.n s (m s z)} die 2en erkennt er jeweils als die Churchzahl 2 und ersetzt sie durch \enquote {$\lambda$s.$\lambda$z.s(s z)}. Dann führt das Programm $\beta$-Reduktion aus und bekommt als Ergebnis\enquote {$\lambda$s.$\lambda$z. (s (s (s (s z))))} dies erkennt er ist die Churchzahl 4 und gibt dies aus. }

\teststep { Der vorherige Test wurde soeben durchgeführt}
{ Der Benutzer drückt auf die drei Punkte um alle Schritte zu sehen die bei der $\beta$-Reduktion durchgeführt wurden }
{ Das Programm zeigt folgende Schritte :
\newline ($\lambda$n.$\lambda$m.$\lambda$s.$\lambda$z.n s (m s z)) ($\lambda$s.$\lambda$z.s(s z)) ($\lambda$s.$\lambda$z.s(s z))
\newline ($\lambda$m.$\lambda$s.$\lambda$z.($\lambda$s.$\lambda$z.s(s z)) s (m s z)) ($\lambda$s.$\lambda$z.s(s z))
\newline $\lambda$s.$\lambda$z.($\lambda$s.$\lambda$z.s(s z)) s (($\lambda$s.$\lambda$z.s(s z)) s z)
\newline $\lambda$s.$\lambda$z.($\lambda$z.s(s z)) (($\lambda$s.$\lambda$z.s(s z)) s z)
\newline $\lambda$s.$\lambda$z. ( s ( s (($\lambda$s.$\lambda$z.s(s z)) s z) ))
\newline $\lambda$s.$\lambda$z. ( s ( s (($\lambda$z.s(s z)) z) ))
\newline $\lambda$s.$\lambda$z. ( s ( s ( s (s z))))
 }

\teststep{Das Programm ist im Default-State}
{der Benutzer gibt in den Editor ($\lambda$x.xx)($\lambda$x.xx) ein }
{ das Programm erkennt die Endlosschleife und warnt den Benutzer }

\teststep {Das Programm ist im Default-State}
{ Der Benutzer drückt auf den \enquote {Plus} Button und gibt \enquote {pair} und $\lambda$a.$\lambda$b.$\lambda$f.f a b ein }
{Wenn er auf den Button drückt öffnet sich ein Popup-Fenster in dem der Benutzer als Name der neuen Funktion Pair eingeben kann und als lambda term den von ihm gewünschten lambda term eingibt. Dieser Name wird mit dem Lambdaterm gespeichert und man kann im Editor anstelle des Lambdaterms den Namen eingeben. }

\teststep { Der vorherige Test wurde ausgeführt}
{ der Benutzer gibt in den Editor \enquote {first (pair a b)} ein. Dannach drückt er auf Run}
{ Das Programm erkennt first aus seiner Standardbibliothek und ihm ist $\lambda$p.p($\lambda$a.$\lambda$b.a) zugewiesen. Das Programm erkennt pair aus dem im letzten Testfall hinzugefügten Lambdaterm und kann nun $\beta$-Reduktion druchführen. Das Programm gibt als Ergebnis a aus.}

\teststep {Das Programm berechnet zurzeit einen Lambdaterm }
{Der Benutzer drückt auf den Stopp-Button}
{ Das Programm beendet die Berechnung und gibt kein Ergebnis aus }

\teststep {Im Editor steht der Term ($\lambda$t.$\lambda$f.f)(($\lambda$y.($\lambda$x.x x))(($\lambda$x.x)($\lambda$x.x)))($\lambda$t.$\lambda$f.f) }
{Der Benutzer benutzt den Step-Button einmal mit Normalreihenfolge  und einmal mit der Auswertungsstrategie Call-By-Value welche in den Einstellungen zu finden ist}
{ Bei der Normalreihenfolge wird der Redex ($\lambda$t.$\lambda$f.f) als erstes reduziert. bei der Auswertungsstrategie wird der Term ($\lambda$x.x) als erstes reduziert }

\teststep { Das Programm ist im Default-State}
{ Der Benutzer gibt ($\lambda$y.$\lambda$x.y x) x a in den Editor ein und drückt auf Run}
{ Durch die Alphakonversion mit De Bruijn Indizes ist das Ergebnis nicht a  a  }

\teststep {Das Programm hat einen Lambda-Term im Editor stehen} { Der Benutzer drückt auf einen Redex}
{ Dieser Redex wird als nächstes ausgewertet unabhängig davon welche Auswertungsstrategie ausgewählt ist. }

\teststep { Das Programm ist im Default-State}
{ Der Benutzer gibt den (($\lambda$x.x)( in den Editor ein }
{ Das Programm gibt die Fehlermeldung aus, dass dies kein gültiger LambdaTerm ist }

\teststep { Das Programm ist im Default -State }
{ Der Benutzer drückt shift l }
{Im Editor erscheint ein $\lambda$}

\teststep {Das Programm hat einen Lambdaterm ausgerechntet}
{ Der Benutzer drückt auf den Share button und wählt copy-link}
{ Eine gekürzte Url wird in sein Clipboard gespeichert }

\teststep {Der Benutzer hat von vorherigem Test eine Url erhalten}
{Der Benutzer gibt den Link in seine URL Leiste ein }
{ Das Programm öffnet sich mit dem LambdaTerm bereits in der Eingabe }

\teststep {Das Programm ist im Default-State}
{ Der Benutzer gibt sub 1 1 ein und drückt auf Run }
{ Das Programm erkennt die funktion substraction und die churchzahlen und gibt als ergebnis 0 aus. Da das Programm nicht erkennen kann ob es die Churchzahl 0 oder der boolean false ist gibt es immer 0 aus. Im Tutorial des Programmes wird das dem Benutzer erläutert }

\teststep {Das Programm ist im Default-State}
{Der Benutzer gibt ($\lambda$x.xx)($\lambda$x.xx) in den Editor ein.}
{Das Programm erkennt die Endlosschleife und terminiert.}

\teststep {Das Programm ist im Default-State}
{Der Benutzer drückt im Feld Benutzerbibliothek auf den Plus-Button.}
{Es öffnet sich ein Popup, indem der Benutzer eine Funktion definieren kann.}




%\teststep{Besucher \enquote{Jayne Cobb} ist auf der Homepage}
%{Er folgt dem Link mit dem Text \enquote{Datenschutz}}
%{Ein Text mit allen Datenschutzinformationen wird ihm angezeigt.}

%\teststep{}
%{Jayne folgt dem Link mit dem Text \enquote{Impressum}}
%{Ein Text mit Informationen des Betreibers wird ihm angezeigt.}

\section{Entwicklungsumgebung}
  \begin{description}
	\item[Benutzte Software die für das Erstellen des Programmes verwendet wurden]~\par
	\begin{itemize}
		\item \textbf{Programmiersprache}
		\begin{itemize}
			\item Java 8
		\end{itemize}
		\item \textbf{Build System}
		\begin{itemize}
			\item Bazel
		\end{itemize}
		\item \textbf{IDE}
		\begin{itemize}
			\item Jeder Entwickler wählt selber. Eclipse, Intellij oder Sublimetext
		\end{itemize}
		\item \textbf{Webframework}
		\begin{itemize}
			\item Google Web Toolkit \enquote{GWT}
		\end{itemize}
		\item \textbf{Versionskontrolle}
		\begin{itemize}
			\item Git auf scc.kit.edu. Webanwendung ist Gitlab
		\end{itemize}
		\item \textbf{Dokument Erstellung}
		\begin{itemize}
			\item LaTeX
		\end{itemize}
		\item \textbf{Testen}
		\begin{itemize}
			\item JUnit
			\item JaCoCo
		\end{itemize}
		\item \textbf{GUI Entwürfe}
			\begin{itemize}
				\item draw.io
			\end{itemize}
		\item \textbf{UMl und Diagramme zeichnen}
			\begin{itemize}
				\item UMLet
				\item Umbrella
				\item BOUML
				
			\end{itemize}
		\item \textbf{endgültige Programmiersprache auf der Website}
			\begin{itemize}
				\item Javascript
			\end{itemize}
	\end{itemize}
%%%%%%%%%%%%%
\newpage
\appendix
\end{description}
\section{Seitenentwürfe}
\textbf{GUI HIER}
% made via https://gomockingbird.com/projects/mnf0cwf/4gXVnC

%\begin{figure}[hb]
%\fbox{\includegraphics[width=\textwidth]{image/login.png}}
%\caption{\label{fig:homepage}
%Homepage mit Login-Funktion
%}
%\end{figure}


\section{Glossar}

%\textbf{Homepage}:
%Seite, die beim Besuchen der Betreiberdomain \emph{ohne Pfad} angezeigt wird. Auch %\enquote{Startseite}.
\textbf{Standardbibliothek}
Eine Ansammlung von Lambdatermen die einen Funktionsnamen bekommen haben. Die Entwickler von Aurora haben die wichtigsten Lambdaterme zusammen getragen damit der Benutzer häufig vorkommende Lambdaterme nicht selber tippen muss. Der Benutzer kann den Namen der Funktion im Editor verwenden und das Programm ersetzt den Namen automatisch in den gewählten Lambdaterm.

\textbf{Benutzerbibliothek}
Der Benutzer kann selber Funktionen mit Namen und Lambdaterm definieren. er kann dann den gewählten Namen im Editor verwenden und der Name wird automatisch durch den Lambdaterm ersetzt.

\textbf {Redex}
Redex steht für \enquote{reducible expression} oder auf deutsch \enquote{reduzibler Ausdruck} ist ein Subterm der durch die Auswertungsstrategien mit Beta-Reductionen reduzierbar sind.
Bei dem Term ($\lambda$x.x) y ist der ($\lambda$x.x) der Redex.

\textbf{Normalform}
Die Normalform ist ein $\lambda$-term der nicht mehr durch Betareduktion reduzierbar ist. Dies wird im Programm als \enquote{Result} ausgegeben

\textbf{Auswertungsstrategie}
Die Auswertungsstrategie gibt an welcher Redex als nächstes ausgewertet wird und dann mit Beta-Reduktion reduziert wird.

\textbf{Normalreihenfolge}
Der linkeste äußerste Redex wird als erstes ausgewertet und reduziert.

\textbf{Call-By-Name}
Der linkeste äußerste Redex der nicht von einem $\lambda$ umgeben ist wird als erstes ausgewertet und reduziert

\textbf{Call-By-Value}
Der linkeste Redex der nicht von einem $\lambda$ umgeben ist und dessen Argument ein Wert ist wird als erstes ausgewertet und reduziert.

\textbf{Hover}
Mit dem Mauszeiger über einem Element \enquote{schewben}, ohne es jedoch zu drücken.
 
\end{document}
